% ===========================================================================
\chapter{Roads as a graph}\label{ch:graph}

To find a way from one place to another, a route planner forgets almost
everything about the map: the colours, the lake, the names of the woods.
What is left is a \emph{graph}.

\section{Graphs}

A graph is a set of \emph{nodes} (also called vertices) joined by
\emph{edges}. For a road map, the nodes are the points along the roads and
the edges are the stretches of road between neighbouring points. An edge
has a direction, from its \emph{source} node to its \emph{target} node: a
two-way road gives two edges, one each way, and a one-way road gives only
the one in the direction of traffic. And every edge has a \emph{cost}:
its length in metres, or the time it takes to travel, which depends on the
kind of road.

\begin{center}
\begin{tikzpicture}[
  n/.style={circle, fill=accent, inner sep=1.8pt},
  e/.style={-{Stealth[length=2.2mm]}, thick, color=cm},
  lab/.style={font=\scriptsize, text=cm}]
  \node[n] (a) at (0,0) {};
  \node[n] (b) at (2.4,0) {};
  \node[n] (c) at (4.8,0) {};
  \node[n] (d) at (2.4,1.8) {};
  \node[n] (e) at (2.4,-1.6) {};
  \draw[e] (a) to[bend left=12] node[above, lab] {100 m, 12 s} (b);
  \draw[e] (b) to[bend left=12] (a);
  \draw[e] (b) to[bend left=12] (c);
  \draw[e] (c) to[bend left=12] (b);
  \draw[e] (b) -- node[right, lab] {one-way} (d);
  \draw[e] (e) to[bend left=12] (b);
  \draw[e] (b) to[bend left=12] (e);
  \node[font=\small, text=cm, anchor=west] at (5.4,0) {a two-way street: two edges};
  \node[font=\small, text=cm, anchor=west] at (2.8,1.8) {a one-way street: one edge};
\end{tikzpicture}
\end{center}

The same map gives different graphs to different travellers. A driver
can't use footpaths and must respect one-way streets. A walker can use
every road and path, in either direction. So the graph is built for a
\emph{mode} of travel: driving or walking.

\section{Speeds}

How fast does a traveller go? Real route planners use the speed limit of
each road, adjusted by the time of day and live traffic. We use one speed
for each kind of road: 60\,km/h on the main road, 40 on the lanes, 30 in
town and 20 on the bumpy tracks, while a walker goes at a steady
5\,km/h anywhere. A speed of 0 means the traveller can't use that kind of
road at all.

\begin{newfile}{road_graph.nex}
-- road_graph.nex: the roads as a graph, for one way of travelling.

intern Island_Map

-- How the traveller gets about.
enum union Mode
  Driving
  Walking
end

-- How fast a traveller goes on a kind of road, in km/h; 0 means the
-- traveller may not use it at all.
function speed_kmh(mode: Mode, kind: Road_Kind): Real do
  if mode = Mode.Walking then
    result := 5.0
  elseif kind = Road_Kind.Main then
    result := 60.0
  elseif kind = Road_Kind.Lane then
    result := 40.0
  elseif kind = Road_Kind.Street then
    result := 30.0
  elseif kind = Road_Kind.Track then
    result := 20.0
  else
    result := 0.0
  end
ensure
  non_negative: result >= 0.0
end
\end{newfile}

An edge knows its two nodes, its length, and the road it belongs to. The
time it takes is worked out once, when the edge is made: a speed of $v$
km/h is $v / 3.6$ metres per second (a kilometre is 1,000 metres and an
hour 3,600 seconds), and time is distance divided by speed.

\begin{append}{road_graph.nex}
-- One step of travel: from node `source` straight to the next node
-- along a road, `target`.
class Edge
create
  make(s, t: Integer, metres, kmh: Real, r: Road)
  require
    different_nodes: s /= t
    positive_length: metres > 0.0
    positive_speed: kmh > 0.0
  do
    source := s
    target := t
    length := metres
    seconds := metres / (kmh / 3.6)
    road := r
  end
feature
  once source: Integer
  once target: Integer
  -- How long the step is, in metres...
  once length: Real
  -- ...and how long it takes, in seconds.
  once seconds: Real
  -- The road the step is part of.
  once road: Road
end
\end{append}

\section{Building the graph}

How should a graph be stored? The question a route planner asks over and
over is ``where can I go from this node?'', so we store, for every node,
the list of edges leaving it. This is called an \emph{adjacency list}. A
map from node ids to arrays of edges does it.

To build the graph, go through every road the traveller may use, and every
pair of neighbouring nodes along it. Add the edge in the road's direction,
and, unless the road is one-way and the traveller is driving, the edge
back.

\begin{append}{road_graph.nex}
class Road_Graph
create
  -- The graph of the roads of `m` that a traveller can use in mode `how`.
  build(m: Island_Map, how: Mode) do
    island := m
    mode := how
    adjacency := {}
    edge_count := 0
    top_speed := 0.0
    across m.roads as r do
      let kmh: Real := speed_kmh(how, r.kind)
      if kmh > 0.0 then
        top_speed := top_speed.max(kmh / 3.6)
        let both_ways: Boolean := not r.is_one_way or how = Mode.Walking
        from
          let i: Integer := 1
        until
          i >= r.node_ids.length
        do
          let a: Integer := r.node_ids.get(i - 1)
          let b: Integer := r.node_ids.get(i)
          let metres: Real := m.point(a).distance_to(m.point(b))
          add_edge(create Edge.make(a, b, metres, kmh, r))
          if both_ways then
            add_edge(create Edge.make(b, a, metres, kmh, r))
          end
          i := i + 1
        end
      end
    end
  end
feature
  island: Island_Map
  mode: Mode
  -- For every node, the edges that leave it.
  adjacency: Map[Integer, Array[Edge]]
  edge_count: Integer
  -- The fastest speed on any edge, in metres per second.
  top_speed: Real
\end{append}

The graph also remembers the fastest speed on any of its edges. That will
be useful in Chapter~\ref{ch:astar}.

The questions the graph answers: is a node on its roads, which nodes it
has, which edges leave a node, and is there an edge from one node to
another. A node with no edges leaving it, such as the last node of a
one-way street that ends at nothing, is still a node of the graph, so
\code{add_edge} makes sure both ends are in the map.

\begin{continued}{road_graph.nex}
  -- Whether `id` is on a road this graph's traveller can use.
  has_node(id: Integer): Boolean do
    result := adjacency.contains_key(id)
  end

  -- Every node of the graph.
  nodes(): Array[Integer] do
    result := adjacency.keys
  end

  -- The edges leaving `id` (none if it is not in the graph).
  edges_from(id: Integer): Array[Edge] do
    result := adjacency.try_get(id, [])
  end

  -- The edge from a to b, if there is one.
  edge_between(a, b: Integer): ?Edge do
    across edges_from(a) as e do
      if e.target = b then
        result := e
      end
    end
  end
\end{continued}

\section{Joining the roads}

A place, or a point someone clicks on the map, is usually not on a road
at all. The harbour is at the water's edge, a little south of Quay
Street; the campsite is in a field at the end of a track. To plan a route
from it, we first join the roads at the nearest node, which is called
\emph{snapping}. The nearest node is simply the one at the smallest
distance; for the few hundred nodes of an island, checking them all is
instant. (A real map has hundreds of millions of nodes, and finds the
nearest with a \emph{spatial index}, which narrows the search to the
nodes in the same small square of the map.)

\begin{continued}{road_graph.nex}
  -- The node of the graph closest to p: where a traveller standing at p
  -- joins the roads.
  nearest_node(p: Point): Integer
  require
    not_empty: edge_count > 0
  do
    let best: Real := -1.0
    across adjacency.keys as id do
      let d: Real := island.point(id).distance_to(p)
      if best < 0.0 or d < best then
        best := d
        result := id
      end
    end
  ensure
    in_graph: has_node(result)
  end

private feature
  add_edge(e: Edge) do
    if not adjacency.contains_key(e.source) then
      adjacency.put(e.source, [])
    end
    -- A one-way road's last node has no edge leaving it, but it is still a
    -- node of the graph.
    if not adjacency.contains_key(e.target) then
      adjacency.put(e.target, [])
    end
    adjacency.get(e.source).add(e)
    edge_count := edge_count + 1
  end
invariant
  non_negative_count: edge_count >= 0
end
\end{continued}

Because the snapping is done in the traveller's own graph, a walker can be
snapped to a footpath, but a driver never is.

\section{Looking at a junction}

Where Harbour Street crosses Market Street, in the middle of town, four
stretches of road meet. Market Street is one-way, eastbound. This script
finds the node the two streets share and lists the edges leaving it, for a
driver and for a walker:

\begin{tryit}{try_graph.nex}
intern Road_Graph
intern Helpers

let island: Island_Map := create Island_Map.load(user_path("kestrel_island.map"))
let driving: Road_Graph := create Road_Graph.build(island, Mode.Driving)
let walking: Road_Graph := create Road_Graph.build(island, Mode.Walking)
console.print_line("driving: " + driving.nodes.length.to_string + " nodes, "
  + driving.edge_count.to_string + " edges")
console.print_line("walking: " + walking.nodes.length.to_string + " nodes, "
  + walking.edge_count.to_string + " edges")

-- The node that Harbour Street and Market Street share.
let crossing: Integer := 0
across island.roads as a do
  across island.roads as b do
    if a.name = "Harbour Street" and b.name = "Market Street" then
      across a.node_ids as id do
        if b.node_ids.contains(id) then
          crossing := id
        end
      end
    end
  end
end
across [driving, walking] as g do
  console.print_line("")
  console.print_line(g.mode.to_string + " from node " + crossing.to_string + ":")
  across g.edges_from(crossing) as e do
    console.print_line("  to node " + e.target.to_string.pad_end(" ", 4) + " along " + e.road.name
      + ", " + bearing(island.point(e.source), island.point(e.target)).round.to_string
      + " degrees, " + e.length.format(0) + " m, " + e.seconds.format(1) + " s")
  end
end
\end{tryit}
\tryoutput{try_graph}

A driver can leave the junction three ways: north and south on Harbour
Street, and east on Market Street. A walker has a fourth: west on Market
Street, against the traffic. The same 100 metres takes 12 seconds to
drive and 72 seconds to walk.

The walking graph has more nodes than the driving one, since it includes
the footpaths, and nearly twice as many edges per road, since every road
is two-way on foot.
